\section{Open questions raised by this replication}

Five concrete follow-up questions surfaced while performing this replication.
Each is stated with its evidentiary basis (why the current work does not
resolve it) and concrete next steps.

\begin{enumerate}

\item \textbf{Non-Pauli noise robustness.} How robust is
Richardson-extrapolation-based ZNE when the underlying noise is genuinely
non-Pauli (amplitude damping, coherent over/under-rotation, leakage,
crosstalk) rather than the symmetric depolarizing channel used here?
\emph{Basis:} our simulator uses symmetric depolarizing noise, which is an
exact polynomial in the stretch factor $c$, so Richardson reaches machine
precision by construction; the paper's own hardware Richardson error is
$\mathcal{O}(10^{-2})$, not $10^{-16}$.
\emph{Next steps:} rerun \texttt{replicate.py} with (a) amplitude damping
plus dephasing, (b) an injected coherent single-qubit over-rotation, and
(c) a two-qubit crosstalk term; measure how each extrapolator degrades and
identify the crossover where linear extrapolation beats Richardson.

\item \textbf{Composition with virtual distillation.} Does composing ZNE
with virtual distillation give a multiplicative bias-reduction advantage,
or do the two techniques compete because they consume overlapping shot
budgets? \emph{Basis:} ZNE cancels bias analytic in the noise rate;
virtual distillation projects onto the dominant eigenvector of $\rho$
giving exponential-in-$M$ suppression. Whether the residual errors are
independent is not settled.
\emph{Next steps:} extend \texttt{replicate.py} to a 2-copy
virtual-distillation estimator of the same $\langle Z_0 Z_1 \rangle$
observable, apply ZNE on top, compare (baseline, ZNE-only, VD-only,
ZNE+VD) at matched total shot budget.

\item \textbf{Chemistry-benchmark scaling.} How does ZNE's practical error
reduction scale on a chemistry-relevant workload (H$_2$ or LiH VQE energy,
or a small transverse-field Ising ground-state energy) once the observable
is a sum of many non-commuting Paulis? \emph{Basis:} our replication tests
a single 2-qubit Pauli on a Bell state; per-term ZNE gains may not
survive summation when variance grows.
\emph{Next steps:} build H$_2$/STO-3G via OpenFermion, run a
hardware-efficient ansatz on a noisy simulator, apply per-term ZNE with a
shared stretch schedule, and compare mitigated vs. unmitigated dissociation
curves; report shot overhead to reach chemical accuracy.

\item \textbf{Principled scale-factor selection.} What is the principled
way to select the noise-scale factors $\{c_1, c_2, \dots\}$ on realistic
hardware where $p_{\mathrm{eff}}$ is not exactly proportional to circuit
depth? \emph{Basis:} both the paper and this replication assume
$p_{\mathrm{eff}}(c) = c\,p_0$; practitioners use ad-hoc grids.
\emph{Next steps:} characterize $p_{\mathrm{eff}}(c)$ empirically via
randomized benchmarking, then formulate scale-factor selection as a
bias-variance-optimal experimental-design problem (D-optimal or
Bayes-optimal) and benchmark against the uniform grid.

\item \textbf{ML-driven extrapolation.} Can a lightweight ML model
(Gaussian process, small MLP, or symbolic regression) trained on
$(c, E(c))$ batches outperform hand-picked Richardson / linear /
exponential extrapolators, especially under non-polynomial noise?
\emph{Basis:} our three extrapolators differ by more than an order of
magnitude; extrapolator choice is currently manual.
\emph{Next steps:} generate synthetic $(c, E(c))$ datasets under a grid of
noise models, train a GP regressor and a shallow MLP to predict $E(0)$
from a fixed-length feature vector, and benchmark against the three
closed-form extrapolators on held-out noise conditions; report win rate
and failure modes (in particular ML overfitting/hallucination).

\end{enumerate}
